% Draft copies only: appended by `make paper1-pdf DRAFT=1` (manuscript.tex, \draftappendix).
% A blind analytic summary (Muller 2017) of this draft, written 2026-11-14 by a language
% model that saw only the extracted text and Muller's handout. Regenerate it for each draft;
% source and inputs: lab/reviews/2026-11-14b-*. Not part of the journal or arXiv build.
\section{Draft only: a blind analytic summary of this draft}
\label{app:analytic}

\noindent\textit{For readers of this draft; it will not appear in the published version.
An analytic summary sets out a paper's claims, arguments, the objections it considers and
where its reasoning is hard to follow (F.A. Muller, \textit{What is an Analytic Summary?}, course handout, 2017). This one was written on 14 November
2026 by a language model that was given only the text of this draft and Muller's description
of the form, with no other information about the paper or its author. Locators (L\ldots) are line numbers of the text extract it was given, not of this PDF. It counts 39 register entries; there are now 41.  Since it was written the draft has changed in several places it comments on, among them the gap in Section~1, the warrant gloss, scope versus objection, the worked example, the sources for the first example, who drafted the register, Table~2's note, and the third error's example (Greenberg's case is now used only for relayed claims).}

\medskip
{\small
\subsection{0. Background}\label{background}

LLM writing assistants are now common in scholarly writing (L17-20). The
literature on their failures has centred on fabricated citations
(L22-27). Journals respond with disclosure rules (L49-50), and the
empirical sciences rely on reporting guidelines such as CONSORT and
PRISMA (L64-68). The paper sits between research-integrity work,
argumentation theory (Toulmin, Johnson and Blair, Walton) and
theory-building methodology (Whetten, Gregor). It borrows the
engineering vocabulary of verification, validation and traceability
(L203-211). It is a perspective paper. It proposes a procedure and a
published artefact, the ``claim register'', and applies the procedure to
itself. It reports no trial with other users.

\subsection{1. Results of conceptual
analyses}\label{results-of-conceptual-analyses}

\begin{enumerate}
\def\labelenumi{\arabic{enumi}.}
\tightlist
\item
  \textbf{Three error types in AI-assisted drafts} (L22-40): (a)
  citation errors; (b) confidence inflation, meaning wording more
  certain than the evidence; (c) overreach, meaning real sources
  correctly attributed but stretched beyond their cases, period or
  setting.
\item
  \textbf{Three kinds of check are kept apart} (L41-49): a citation
  check (does the work exist and say this?), a reporting guideline (is
  the study fully reported?), and a reasoning check (does the step from
  evidence to claim hold?). Authorship detection is a fourth, separate
  question (who drafted it?).
\item
  \textbf{Non-empirical paper} (L79-81): a paper whose main contribution
  is reasoning (theoretical, conceptual, design-science, perspective).
  The concept is motivated by Gregor's Types I and V. Such a paper's
  whole is new by definition, so no earlier source states it.
\item
  \textbf{Toulmin adapted} (L102-111, Box 1): a claim rests on grounds
  through a warrant, with a qualifier and a scope. ``Scope'' renames
  Toulmin's rebuttal. ``Grounds'' extends Toulmin's data to four kinds:
  source, other claims, declared starting point, own work. Backing is
  recorded only when a warrant may be disputed.
\item
  \textbf{Warrant check} (L154-165): relevance (each ground bears on the
  claim) and sufficiency (all grounds together justify the claim at its
  stated strength). These come from Johnson and Blair's notion of
  cogency. The check has four outcomes: pass, pass confirmed by another,
  named gap, fail (L166-169). Failure means ``grounds do not carry the
  claim as worded'', not ``claim is false''.
\item
  \textbf{Confidence tier} (L180-183, Table 1): the lower of two
  ratings, support for the grounds and the outcome of the warrant check.
  Four tiers (Established, Supported, Emerging, Speculative), each
  licensing fixed qualifiers.
\item
  \textbf{Priority} (L190-193): P0/P1/P2 by what fails if the claim is
  wrong. Every P0 entry should reach Supported.
\item
  \textbf{Verification versus validation} (L203-207): checking each
  claim is verification; checking the contribution with Whetten's
  questions is ``closer to validation''. The register is thus a
  traceability record.
\end{enumerate}

\subsection{2. Claims}\label{claims}

\textbf{Central thesis (T0).} A claim register published with a paper
whose contribution is an argument can make that paper's reasoning
traceable and checkable step by step. Registers should be trialled
before journals adopt them (L53-58, L300-308; Table 3, C1).

\begin{itemize}
\tightlist
\item
  \textbf{T1 (gap).} No consensus checklist exists for checking the
  reasoning of non-empirical papers step by step (L51-52, L93-98; S2-2).
\item
  \textbf{T2.} Citation checks and reporting guidelines do not check the
  step from evidence to claim (L41-44, L69-75).
\item
  \textbf{T3.} In an argument paper, an unchecked reasoning step means
  an unchecked contribution (L42-44).
\item
  \textbf{T4.} AI detection does not fill the gap (L45-49).
\item
  \textbf{T5.} AI-assisted drafts show citation errors, possibly
  confidence inflation, and overreach (L22-40; S1-1).
\item
  \textbf{T6.} A check chosen by kind of grounds may avoid both false
  passes and false failures that citation checks or source counts
  produce (L212-217; S3-4).
\item
  \textbf{T7 (novelty).} Each of the register's four parts is used
  elsewhere, but none of the compared approaches combines them for
  academic papers (L219-221, L265-266; Table 2).
\item
  \textbf{T8.} A passed register shows method was followed, not that
  claims are true. Established needs confirmation by someone else
  (L185-189).
\item
  \textbf{T9 (self-application).} On its own register, the paper's three
  central P0 entries stay at Emerging (L271-275).
\item
  \textbf{T10.} AI-drafted registers risk replacing rather than exposing
  the author's reasoning (L289-298).
\end{itemize}

\subsection{3. Arguments}\label{arguments}

\textbf{For T1.} Empirical premise: a dated search of EQUATOR (709
titles) and four databases, 28 queries, 701 titles screened by one
person, found none (L93-98). Supporting premise: scattered statements in
conceptual, design-science and software-engineering methodology that no
agreed guidance exists (L85-92). {[}RECONSTRUCTED{]} An absence of hits
in these sources makes absence in the literature likely. The paper
limits this itself: one screener, no Scopus, Web of Science or
humanities indexes (L96-98), so it rates T1 Emerging (Table 3).

\textbf{For T2.} Definitional premise: reporting guidelines guide
reporting (L69-71, Moher). Empirical premise: CONSORT's measured effect
is on completeness (L71-73). Conclusion: they ``do not ask whether the
reasoning holds'' (L74-75). For citation checks the argument is the
worked example (L112-128): each citation passes, yet the claim
overreaches.

\textbf{For T3.} In an argument paper the step from evidence to claim is
``what the paper adds'' (L43-44). {[}RECONSTRUCTED{]} If no instrument
checks that step, nobody checks it. This needs the further premise that
peer review does not check reasoning. The paper later denies that
premise (L314).

\textbf{For T4.} A detector answers who drafted a text, not whether it
holds (L48-49). The argument is a category distinction and does not
depend on detector accuracy.

\textbf{For T5.} Citation errors: two studies (L23-27). Confidence
inflation: LLMs verbalise high confidence on wrong answers (L29-31), but
evidence for academic prose is mixed (L31-33). Overreach: a biomedical
citation-distortion case (L39-40) and the worked example.
{[}RECONSTRUCTED{]} The text needs, but does not state, that AI
assistance makes overreach more frequent or harder to see. Its only
support is that AI drafts ``read as polished'' (L20-21).

\textbf{For T6.} Premise: a citation check or a source count does not
inspect the step (L212-214). Premise: a check matched to the grounds
(source against scope, inference against qualifier) does inspect it
(L214-216). Illustration: the worked example as a false pass caught
(L216-217). The argument is hedged (``may''). It gives no example of a
false failure avoided.

\textbf{For T7.} A feature comparison of eight approaches across four
parts (Table 2, L222-264). Each part has users, and no row except the
register's has all four. {[}RECONSTRUCTED{]} The comparison set covers
the relevant prior art. The scope phrase ``compared here'' limits the
claim to this set, but L266 drops the hedge.

\textbf{For T0.} Two components. (i) The capability (``can make
reasoning traceable''): by construction, since the register records each
link (L207-211), plus the engineering analogy (NASA, IPCC ``traceable
account''). (ii) The normative part (``trial first''): from the declared
value that a proposal be tested before adoption (D7, Table 4) and from
T9 and the limitations (L277-280). C1 rests on five claims and six
declared values and definitions (Table 3). Its tier follows a
weakest-link rule {[}RECONSTRUCTED from ``no higher than its weakest
supporting claim'', Table 3{]}.

\textbf{For T8.} Self-checks can err in the author's favour (L187-188).
Therefore Established requires third-party confirmation, normally peer
review (L188-189).

\textbf{For T10.} Writing a warrant forces the author to state the step;
checking a tool-written warrant asks less (L292-294). Automation
complacency findings (Parasuraman and Manzey) and self-reported reduced
critical thinking (Lee et al.) support this (L295-298).

\subsection{4. Objections considered and
responses}\label{objections-considered-and-responses}

\begin{longtable}[]{@{}
  >{\raggedright\arraybackslash}p{(\linewidth - 4\tabcolsep) * \real{0.3333}}
  >{\raggedright\arraybackslash}p{(\linewidth - 4\tabcolsep) * \real{0.3333}}
  >{\raggedright\arraybackslash}p{(\linewidth - 4\tabcolsep) * \real{0.3333}}@{}}
\toprule\noalign{}
\begin{minipage}[b]{\linewidth}\raggedright
Objection
\end{minipage} & \begin{minipage}[b]{\linewidth}\raggedright
Author's response
\end{minipage} & \begin{minipage}[b]{\linewidth}\raggedright
Assessment
\end{minipage} \\
\midrule\noalign{}
\endhead
\bottomrule\noalign{}
\endlastfoot
AI detection already addresses AI writing (L45-49) & Detection answers
who drafted, not whether it holds & Adequate \\
Reporting guidelines and field guidance already cover this (L64-92) &
They target completeness; reasoning guidance is scattered and lacks
templates & Adequate for reporting guidelines; Partial for JBI and
SANRA, which do judge reasoning (L231-237) \\
Toulmin did not intend his layout as an analytical model (L110-111) &
Using it as a checklist is ``an adaptation'' & Weak: the adaptation is
named but not defended \\
Fixed qualifiers may not secure a shared reading (Budescu, L176-179) &
The table ``may help an author notice'' overclaiming (L183-184) &
Partial: it moves the function from readers to authors without saying
so \\
A passed register does not make claims true (L185-188) & Conceded;
Established needs outside confirmation & Conceded \\
The author checks his own rules: circularity (L286-288) & The warrant
check's reasons are recorded so others can disagree & Partial: this
enables later correction, not independence now \\
Existing approaches already do this (Section 4) & Novelty lies in the
combination (L265-266) & Partial: it depends on the chosen comparison
set and on ``step by step'' \\
Only one user; not tested (L277-278) & Trials recommended before
adoption (L303-308) & Conceded \\
LLM reviewers are not independent of the method (L279-280) & Stated as a
limitation & Conceded \\
Burden on authors: 39 entries (L281-282) & A shorter reviewer form may
suit many papers & Partial: the short form drops the published register,
which is the proposal \\
Western frameworks; poor fit for humanities (L283-285) & Fit ``needs to
be assessed'' & Open \\
AI tools can draft registers fluently (L289-298; Table 4) & Authors are
responsible; source checks need the source in hand & Weak: the cited
complacency evidence holds ``even when told to verify'' (L295-296),
which undercuts responsibility as a safeguard \\
Peer review already judges arguments (L314; Table 4) & The register
supports review and does not replace it & Partial: the response is
consistent, but it conflicts with T3 (see 5.2) \\
Whetten's questions fit a recommendation poorly (L199-202) & Add the
value behind it and its conditions of application & Partial: an ad hoc
extension with no stated criterion \\
Confidence inflation in academic prose is unsettled (L31-33) & None; it
is still listed as an error type & Open \\
\end{longtable}

\subsection{5. Reader's difficulties}\label{readers-difficulties}

\begin{enumerate}
\def\labelenumi{\arabic{enumi}.}
\tightlist
\item
  \textbf{The gap narrows to a qualifier.} T1 holds only for
  ``consensus'', ``non-empirical'' and ``step by step'' together. JBI is
  consensus-based, asks whether an argument is logically defended, and
  points to Toulmin (L231-234). Kane checks step by step (Table 2).
  VANRA verifies assertions in narrative reviews (L225-227). Remove any
  one qualifier and the gap closes. The title's ``gap'' therefore rests
  on a narrow absence, and the text never argues why step-by-step
  checking matters more than whole-argument appraisal.
\item
  \textbf{T3 versus peer review.} L44 says that if no one checks the
  reasoning, no one checks the contribution. That implies nobody does.
  Yet L314 grants that ``peer review already judges arguments'', and
  L189 makes peer review the place of confirmation. The gap is therefore
  not an absence of checking. It is an absence of a structured,
  published record. The paper never states the thesis in that weaker
  form.
\item
  \textbf{Source check versus warrant check in the worked example.} The
  source check already asks whether a source says what is attributed to
  it ``for the cases and period the claim covers'' (L138-140). That
  alone catches the overreach in the example. Yet the example is argued
  in warrant-check terms (relevance, sufficiency; L119-123). The warrant
  check is reserved for steps from other claims or from several sources
  ``to a point none of them makes alone'' (L152-153), which does not
  describe the example. So the example does not show the warrant check
  doing distinctive work. It also weakens T2: a thorough citation check,
  as the paper itself defines one (L41), would also catch overreach.
\item
  \textbf{The warrant gloss is narrow.} Box 1 glosses a warrant as
  ``what holds for these cases holds for this one''. That is
  general-to-particular. The worked example's warrant runs
  particular-to-general (``what the models tested did, LLMs in general
  do'', L116). Analogical, causal and normative warrants, which a
  recommendation paper needs, are not covered (cf.~L163-165).
\item
  \textbf{``Rebuttal'' is split in two.} Box 1 equates Toulmin's
  rebuttal with scope. Figure 1 and L162 add a separate ``strongest
  objection with its answer''. This is the dialectical sense of
  rebuttal, not Toulmin's exception clause. The text does not mark the
  shift.
\item
  \textbf{The tier rule is underspecified for argued claims.} Table 1
  rates ``support for the grounds'' only by counting sources. L181-183
  says an argued claim can reach Supported when its premises are
  supported, but no row says how claim-grounds or declared starting
  points are rated. C1's tier rests on a weakest-link rule that appears
  only in Table 3. Table 1's note lets a one-source claim reach
  Established, which contradicts the ``three or more'' row.
\item
  \textbf{The paper fails its own floor, yet asserts.} L192-193 says a
  paper whose central claims stay at Emerging falls short. L271-273
  reports that this paper's gap claim, central argument and contribution
  do. Still, L266 (``The claim register does'') and Key Point L12-13
  state T7 without a qualifier, and T7's tier is never reported. By the
  paper's own rule, the bare assertion is unlicensed unless S4-9 is
  Supported. The text does not say.
\item
  \textbf{Traceability is not delivered in the text.} Table 4 cites
  S1-7, S2-5, S3-4, S3-6, S4-6 to S4-9, D1 and D7, which the paper never
  shows. Of 39 entries, three appear. ``Three of six P0 reach
  Supported'' (L271) cannot be checked from the paper alone. A reader
  must go to the repository, which the paper itself treats as the
  publication route (L268-269).
\item
  \textbf{The AI link is motivational, not argued.} The overreach error
  is ``not new'' (L39). The register is not AI-specific. The search
  (L93-98) did not concern AI. No premise shows that AI raises the rate
  or lowers the visibility of reasoning errors, beyond fluency (L20-21)
  {[}see 3, T5{]}.
\item
  \textbf{Self-application compounds the T10 risk.} The register was
  ``drafted'' by an LLM (L323-324), which is the scenario T10 warns
  about. Its reviewers were also LLMs (L279, L326-328). The paper does
  not say who wrote the warrants, which by its own analysis decides
  whether the register exposed or replaced the reasoning (L292-294).
\item
  \textbf{Structural displacement.} L129-135 (``Others have put his
  elements to work'') follows the worked example. ``His'' has no
  antecedent there (Toulmin is last named at L110). The paragraph
  previews Section 4 in the middle of Section 3.
\item
  \textbf{``Shows'' in the author's own voice.} L185 (``shows that the
  paper followed the method'') uses a qualifier that Table 1 reserves
  for Established. Read as a definitional claim this is defensible, but
  the text does not mark it as one.
\end{enumerate}

\subsection{6. Title check}\label{title-check}

\begin{itemize}
\tightlist
\item
  \textbf{``The Verification Gap.''} Delivered in part. Section 2
  documents a search and finds no consensus step-by-step checklist
  (L93-98). The paper rates this claim only Emerging (Table 3). The gap
  is defined by narrow qualifiers (5.1). Its stronger framing at L44,
  that no one checks the reasoning, conflicts with L314.
  ``Verification'' acquires its technical sense only at L203-207. Yet
  the gap in the title also includes validation-level concerns (the
  contribution check), which the title leaves out.
\item
  \textbf{``Checking the Reasoning.''} Delivered as a procedure, not as
  a demonstrated check. Section 3 specifies the checks, tiers and
  priorities. The core reasoning check is a recorded judgement on two
  criteria (L154-162, L286-287). Pattern-specific questions are ``not
  yet worked out'' (L164-165). The contribution check fits
  recommendation papers poorly (L199-202). No evidence shows that the
  check finds weak steps (L277-278). The single worked example shows a
  scope error that a source check would also catch (5.3).
\item
  \textbf{``AI-Assisted.''} Weakly delivered. AI frames the introduction
  (L17-50) and returns in the limitations (L289-298) and the
  acknowledgement. The register, the search and the comparison are
  AI-neutral. Two of the three error types are not shown to be
  AI-specific (L31-33, L39). Section 4 covers AI-writing checkers in one
  row. The body supports a broader title about argument papers in
  general.
\item
  \textbf{``Academic Writing.''} Over-broad. The scope is ``scholarly
  articles whose contribution is an argument'' (L59-60), defined as
  non-empirical (L79-81). Empirical papers are mentioned only as
  possible extensions (L312-313). Humanities are flagged as a poor fit
  (L283-285) and were not searched (L97). Student essays and journalism
  appear as speculation (L61-62, L312-313). ``Academic writing'' thus
  names a wider domain than the one argued for.
\end{itemize}

\subsection{7. Audience}\label{audience}

The text names four audiences, with shifting weight.

\begin{itemize}
\tightlist
\item
  \textbf{Authors} are the primary agents: they ``keep'' the register
  (L14, L55), fill it in and are responsible for each entry (L289).
  Table 1 is said to help ``an author notice'' overclaiming (L183-184).
\item
  \textbf{Reviewers and readers} are named as users who ``follow each
  claim back'' (L14-15, L56, L301). The trials, however, test only
  authors and reviewers (L305-306). Readers are never tested, and the
  paper says nothing about what a reader does with 39 entries. The
  shorter form is a reviewer tool that needs no register (L281-282).
  This makes the reviewer a separate audience with a separate artefact.
\item
  \textbf{Journals and editors} ``decide whether to ask for one'' (L57).
  They are addressed in the recommendation to publish registers as
  optional supplements and, conditionally, to build EQUATOR-like
  infrastructure (L303-308). The comparison with reporting guidelines
  (Section 2) is pitched to this audience.
\item
  \textbf{Others.} Teachers and students (L313-314), journalism (L61-62,
  L312-313) and methodologists who would ``test and revise'' the
  candidate (L58-59, L317-319). The last group is arguably the real
  addressee of a paper that asks for trials.
\end{itemize}

\textbf{Consistency.} The text is not fully consistent. Box 1 and L53
define the register as ``published with the paper''. The recommendation
makes it optional (L304). The short reviewer form drops publication
altogether (L281-282). The key points put authors first and treat
journals as an afterthought (L14-15). The recommendations put journals
and trial designers first (L303-308). The AI framing of the title
implies authors who use AI. The procedure is written for all authors of
argument papers. And the acknowledgement shows the register itself was
drafted by an AI (L323-324), so the ``author'' who keeps it is in
practice an author-plus-tool, a role the audience statements never name.
}
